\section{Función objetivo de Flow Matching}

\subsection{Método ingenuo: entrenamiento con máxima verosimilitud global}

Una intuición es maximizar directamente la verosimilitud logarítmica de los datos:

$$\max_\theta \; \mathbb{E}_{x_1 \sim p_{\text{data}}} \left[ \log p_1(x_1) \right]$$

En concreto:
\begin{enumerate}
    \item Se parametriza el campo vectorial con una red neuronal $v_\theta(t, x)$
    \item Para cada punto de datos $x_1$, se ejecuta la ODE hacia atrás desde $x_1$ hasta $x_0$
    \item Se calcula $\log p_0(x_0) - \int_0^1 \nabla \cdot v_\theta(t, \phi_t) \, dt$
    \item Se maximiza este valor
\end{enumerate}

\begin{warningbox}{Defecto letal del método ingenuo}
Aunque este método es correcto en teoría, presenta problemas graves en la práctica:
\begin{enumerate}
    \item \textbf{Coste computacional elevado}: calcular la función de pérdida requiere resolver una ODE en cada paso, lo cual es una operación costosa por sí misma.
    \item \textbf{Inestabilidad numérica}: los solucionadores de ODE tienden a acumular errores numéricos durante integraciones prolongadas.
    \item \textbf{Dificultad para la retropropagación}: el gradiente debe propagarse hacia atrás a través de todo el proceso de resolución de la ODE (requiere el método adjunto), aumentando aún más el coste computacional y la inestabilidad.
\end{enumerate}
Esta es la dificultad central que enfrentaron los primeros Continuous Normalizing Flows (como Neural ODE, Chen et al., 2018).
\end{warningbox}

\subsection{Objetivo de Flow Matching: regresión del campo vectorial}

El avance clave de Flow Matching radica en: \textbf{saltar la resolución de la ODE}, entrenando el campo vectorial directamente con una función de pérdida simple de regresión.

La función de pérdida de Flow Matching se define como:

$$\mathcal{L}_{\text{FM}}(\theta) = \mathbb{E}_{t \sim \mathcal{U}[0,1], \, x_t \sim p_t} \left\| v_\theta(t, x_t) - v_t(x_t) \right\|^2$$

Donde los símbolos tienen el siguiente significado:
\begin{itemize}
    \item $v_\theta(t, x_t)$: campo vectorial predicho por la red neuronal
    \item $v_t(x_t)$: campo vectorial objetivo (ground truth)
    \item $t \sim \mathcal{U}[0,1]$: muestreo uniforme de pasos de tiempo
    \item $x_t \sim p_t$: muestreo de la distribución marginal en el instante $t$
\end{itemize}

\begin{importantbox}{Ventaja clave de Flow Matching}
Esta función de pérdida es una simple \textbf{regresión de error cuadrático medio}, sin necesidad de resolver ninguna ODE. Sin embargo, el problema radica en que no conocemos la forma analítica del campo vectorial marginal $v_t(x_t)$ ni de la distribución marginal $p_t(x_t)$, ya que la propia distribución de datos $p_{\text{data}}$ es desconocida.

Esto da lugar a la innovación clave de Flow Matching: el \textbf{Conditional Flow Matching}.
\end{importantbox}

\subsection{Resumen de la sección}

Maximizar directamente la verosimilitud requiere costosas resoluciones de ODE, mientras que Flow Matching propone una función objetivo de regresión concisa. No obstante, la no analiticidad de las magnitudes marginales exige introducir aún más una formulación condicional.

